\section{Группы Галуа как обратные пределы}

\begin{definition}[7.20]
Частичный порядок $\le$ на множестве $I$ \emph{направленный} (а $(I,\le)$~--- \emph{направленное множество}), если для
любых $i,j\in I$ найдётся $k\in I$ с $i,j\le k$.
\end{definition}

\begin{definition}[7.21]
Пусть $(I,\le)$~--- направленное множество и $\mathsf C$~--- категория (например, групп и гомоморфизмов или
топологических групп и непрерывных гомоморфизмов).
\begin{enumerate}[label=(\alph*),leftmargin=*]
\item \emph{Обратная система} в $\mathsf C$ с индексами $(I,\le)$~--- это семейство объектов $(A_i)_{i\in I}$ и морфизмов
$p^j_i\colon A_j\to A_i$ ($i\le j$) таких, что $p^i_i=\id_{A_i}$ и $p^j_i\circ p^k_j=p^k_i$ при $i\le j\le k$.
\item Объект $A$ вместе с морфизмами $p_j\colon A\to A_j$, для которых $p^j_i\circ p_j=p_i$ при $i\le j$, называется
\emph{обратным} (или \emph{проективным}) \emph{пределом} системы, если он обладает универсальным свойством: для любого
объекта $B$ и морфизмов $q_j\colon B\to A_j$ с $p^j_i\circ q_j=q_i$ существует единственный морфизм $r\colon B\to A$
с $p_j\circ r=q_j$ для всех $j$:
\[
\begin{tikzcd}
B\arrow[r,dashed,"r"]\arrow[rd,"q_j"']\arrow[rdd,bend right=25,"q_i"'] & A\arrow[d,"p_j"]\arrow[dd,bend left=40,"p_i"]\\
 & A_j\arrow[d,"p^j_i"]\\
 & A_i
\end{tikzcd}
\]
\end{enumerate}
Предел (если существует) определён однозначно с точностью до единственного изоморфизма; его обозначают
$\varlim(A_i,p^j_i)$ или $\varlim A_i$.
\end{definition}

\begin{example}[7.22]
Пусть $(G_i,p^j_i)$~--- обратная система групп. Положим
\[
G=\{(g_i)\in\textstyle\prod_iG_i\mid p^j_i(g_j)=g_i\ \text{при }i\le j\}
\]
(\emph{согласованные семейства}) и $p_i\colon G\to G_i$~--- проекции. Равенство $p^j_i\circ p_j=p_i$~--- это условие
$p^j_i(g_j)=g_i$. Если $(H,q_i)$~--- другое такое семейство, то образ гомоморфизма $h\mapsto(q_i(h))\colon
H\to\prod G_i$ лежит в $G$, и это единственный гомоморфизм, переводящий $q_i$ в $p_i$. Значит, $(G,p_i)=\varlim G_i$.
\end{example}

\begin{example}[7.23]
Пусть $(G_i,p^j_i)$~--- обратная система топологических групп и непрерывных гомоморфизмов. С топологией произведения
$\prod G_i$~--- топологическая группа, а $G$ из 7.22 с топологией подпространства~--- её подгруппа; проекции $p_i$
непрерывны. Если $(H,q_i)$~--- другое семейство, то $h\mapsto(q_i(h))$ непрерывно (непрерывны все его компоненты),
значит, $H\to G$ непрерывно, и $(G,p_i)=\varlim G_i$ в категории топологических групп.

Обратную систему конечных групп рассматривают как систему топологических групп, снабжая каждую конечную группу
дискретной топологией. Тогда $G=\varlim G_i$ \emph{замкнуто} в $\prod G_i$: если $(x_i)\notin G$, скажем
$p^{j_0}_{i_0}(x_{j_0})\ne x_{i_0}$, то открытая окрестность $\{(g_i)\mid g_{j_0}=x_{j_0},\ g_{i_0}=x_{i_0}\}$ точки
$(x_i)$ не пересекается с $G$.
\end{example}

\begin{definition}[7.24]
\emph{Проконечная группа}~--- топологическая группа, изоморфная обратному пределу конечных групп (с дискретной
топологией). (<<Проконечная>>~--- сокращение от <<проективный предел конечных>>.)
\end{definition}

\begin{lemma}[7.25a]
В компактной хаусдорфовой вполне несвязной группе $G$ открытые подгруппы образуют базу окрестностей единицы. В частности,
пересечение всех открытых подгрупп равно $\{e\}$.
\end{lemma}

\begin{proof}
\emph{Шаг 1 (топология): открыто-замкнутые множества образуют базу.} Пусть $x\in G$ и $C$~--- пересечение всех
открыто-замкнутых множеств, содержащих $x$. Покажем, что $C$ связно. Пусть $C=A\sqcup B$, где $A,B$ замкнуты в $C$ (а значит, и в $G$, так как $C$ замкнуто), $x\in A$.
Компактное хаусдорфово пространство нормально, поэтому есть непересекающиеся открытые $U\supset A$, $W\supset B$. Так как
$C\subset U\cup W$, а дополнение к $U\cup W$ компактно и покрыто дополнениями открыто-замкнутых множеств, содержащих $x$,
уже конечное пересечение $K$ таких множеств лежит в $U\cup W$; $K$ открыто-замкнуто и содержит $x$. Множество
$K\cap U=K\setminus W$ открыто и замкнуто и содержит $x$, поэтому $C\subset K\cap U\subset U$, откуда $B\subset U\cap
W=\varnothing$. Итак, $C$ связно, и по вполне несвязности $C=\{x\}$. Если $O\ni x$ открыто, то тем же рассуждением с
компактностью конечное (т.\,е.\ открыто-замкнутое) пересечение уже лежит в $O$.

\emph{Шаг 2 (группа): внутри компактной открытой окрестности $K$ единицы есть открытая подгруппа.} Для каждого $k\in K$ по
непрерывности умножения есть открытые $A_k\ni k$ и $W_k\ni e$ с $A_kW_k\subset K$. Конечное число $A_k$ покрывает $K$;
пусть $W$~--- пересечение соответствующих $W_k\cap W_k^{-1}$. Это симметричная открытая окрестность $e$ с $KW\subset K$.
Тогда $KW^n\subset K$ для всех $n$, а так как $e\in K$, то $W^n\subset K$. Подгруппа $H=\bigcup_nW^n$ открыта и лежит в
$K$.

Вместе: любая окрестность $e$ содержит открыто-замкнутую (компактную) окрестность, а та~--- открытую подгруппу.
\end{proof}

\begin{proposition}[7.25]
Топологическая группа проконечна тогда и только тогда, когда она компактна, хаусдорфова и вполне несвязна.
\end{proposition}

\begin{proof}
\emph{$\Rightarrow$.} Пусть $G=\varlim G_i$, $G_i$ конечны. По 7.23 $G$ замкнута в произведении, которое компактно
(Тихонов) и хаусдорфово; значит, $G$ компактна и хаусдорфова. Ядро $U_i$ непрерывной проекции $G\to G_i$~---
открытая подгруппа, следовательно, и замкнутая (7.1a), и $\bigcap U_i=\{e\}$. Как в доказательстве 7.8, компонента
связности единицы лежит в каждой $U_i$, т.\,е.\ равна $\{e\}$; сдвигами~--- то же для любой точки.

\emph{$\Leftarrow$.} Пусть $G$ компактна, хаусдорфова и вполне несвязна.
\begin{itemize}[leftmargin=*]
\item По лемме 7.25a открытые подгруппы $U$ образуют базу окрестностей $e$. Каждая $U$ имеет конечный индекс (7.1a(d)).
\item Сопряжённые к $U$ подгруппы $gUg^{-1}$ зависят лишь от смежного класса $gU$ ($guUu^{-1}g^{-1}=gUg^{-1}$),
поэтому их не больше $(G:U)$. Их пересечение $V=\bigcap_ggUg^{-1}$~--- ядро действия $G$ сдвигами на конечном множестве
$G/U$; значит, $V$~--- нормальная подгруппа конечного индекса ($G/V\hookrightarrow\Sym(G/U)$), открытая как конечное
пересечение открытых, и $V\subset U$. Итак, открытые нормальные подгруппы $V$ тоже образуют базу окрестностей $e$, и
$\bigcap V=\{e\}$.
\item Открытые нормальные подгруппы упорядочены обратным включением и образуют направленное множество
($V_1\cap V_2$~--- снова открытая нормальная). Рассмотрим $r\colon G\to\varlim G/V$, $g\mapsto(gV)_V$ (с отображениями
факторизации $G/V'\to G/V$ при $V'\subset V$). Он непрерывен (каждая проекция $G\to G/V$ непрерывна: $V$ открыта) и
инъективен ($\bigcap V=\{e\}$).
\item \emph{Образ плотен.} Базисная окрестность точки $x=(x_V)$ предела задаётся конечным числом координат:
$\{y\mid y_{V_1}=x_{V_1},\dots,y_{V_k}=x_{V_k}\}$. Положим $W=V_1\cap\dots\cap V_k$ (снова открытая нормальная) и выберем
$g$ с $gW=x_W$. По согласованности $x_{V_i}$~--- образ $x_W$ в $G/V_i$, т.\,е.\ $x_{V_i}=gV_i$. Значит, $r(g)$ лежит в
окрестности.
\item Образ компактен, значит, замкнут; замкнутое плотное множество~--- всё пространство. Непрерывная биекция компакта на
хаусдорфово пространство~--- гомеоморфизм. Итак, $G\cong\varlim G/V$.\qedhere
\end{itemize}
\end{proof}

\begin{remark}
Лемма 7.25a~--- единственное место, где используется нетривиальный факт общей топологии. Компактность в ней существенна:
группа $\Q$ с обычной топологией вполне несвязна, но любая её открытая подгруппа содержит интервал и потому совпадает с
$\Q$; пересечение открытых подгрупп равно $\Q\ne\{0\}$. (Для локально компактных групп верно более общее утверждение:
компонента единицы равна пересечению открытых подгрупп.)
\end{remark}

\begin{example}[7.26: группа Галуа как обратный предел]
Пусть $\Om/F$~--- расширение Галуа. Композит двух конечных подрасширений Галуа~--- снова конечное расширение Галуа (3.21),
поэтому конечные подрасширения Галуа $E\subset\Om$, упорядоченные включением, образуют направленное множество $I$. Для
$E\subset E'$ ограничение $p^{E'}_E\colon\Gal(E'/F)\to\Gal(E/F)$ превращает $(\Gal(E/F))_{E\in I}$ в обратную систему
конечных групп. Ограничения $p_E\colon\Gal(\Om/F)\to\Gal(E/F)$ согласованы ($p^{E'}_E\circ p_{E'}=p_E$) и по
универсальному свойству задают гомоморфизм
\[
r\colon\Gal(\Om/F)\to\varlim_{E\in I}\Gal(E/F),\qquad\sigma\mapsto(\sigma|_E)_E .
\]
\emph{Он является изоморфизмом топологических групп.} Это переформулировка доказательства 7.8 (там вместо $E$ стоят
инвариантные множества $S$, а $G/G(S)\cong\Gal(F(S)/F)$, см.\ 7.8a(a)); проверим и непосредственно.
\begin{itemize}[leftmargin=*]
\item \emph{Инъективность:} каждый $a\in\Om$ лежит в некотором $E\in I$.
\item \emph{Непрерывность:} все $p_E$ непрерывны (7.7).
\item \emph{Плотность образа.} Пусть $x=(x_E)$~--- точка предела и $U=\{y\mid y_{E_i}=x_{E_i},\ i\le k\}$~--- её базисная
окрестность. Положим $E'=E_1\cdots E_k\in I$. По согласованности $x_{E_i}=x_{E'}|_{E_i}$. По 7.7 найдётся $\sigma\in G$
с $\sigma|_{E'}=x_{E'}$; тогда $\sigma|_{E_i}=x_{E_i}$, т.\,е.\ $r(\sigma)\in U$.
\item Образ компактен, значит, замкнут, и, будучи плотным, совпадает с пределом; $r$~--- гомеоморфизм по компактности.
\end{itemize}
Отличие от 7.25: там сюръективность $G\to G/V$ тривиальна, а здесь сюръективность $G\to\Gal(E'/F)$ (7.7) опирается на
7.4~--- продолжение автоморфизмов на $\Om$ по лемме Цорна. В этом теоретико-полевое содержание примера: \emph{элемент
группы Галуа~--- это согласованный набор автоморфизмов всех конечных подрасширений.}
\end{example}

\begin{lemma}[7.26a: конфинальность]
Пусть $\Om=\bigcup_kE_k$, где $E_1\subset E_2\subset\cdots$~--- конечные подрасширения Галуа. Тогда
$\sigma\mapsto(\sigma|_{E_k})_k$~--- изоморфизм топологических групп $\Gal(\Om/F)\cong\varlim_k\Gal(E_k/F)$, а
подгруппы $\Gal(\Om/E_k)$ образуют базу окрестностей единицы. То же верно для любого семейства конечных подрасширений
Галуа, \emph{конфинального} в $I$ (каждое $E\in I$ содержится в одном из полей семейства).
\end{lemma}

\begin{proof}
Каждое конечное подрасширение порождено конечным числом элементов и потому лежит в некотором $E_k$; значит,
$\Gal(\Om/E_k)\subset\Gal(\Om/E)$, и $\{\Gal(\Om/E_k)\}$~--- база. Доказательство 7.26 дословно проходит для системы
$(\Gal(E_k/F))_k$: в шаге о плотности вместо композита берётся наибольшее из $E_{k_1},\dots,E_{k_m}$.
\end{proof}

\begin{remark}[7.26b]
Если $\Om/F$ бесконечно, то $\Gal(\Om/F)$ несчётна. Действительно, выберем строго возрастающую цепочку конечных
подрасширений Галуа $E_1\subsetneq E_2\subsetneq\cdots$ ($E_1=F$; если $E_k$ построено, возьмём $a\in\Om\setminus E_k$ и положим
$E_{k+1}$ равным композиту $E_k$ с полем разложения минимального многочлена $a$) и $\Om'=\bigcup E_k$. Ограничение
$\Gal(\Om/F)\to\Gal(\Om'/F)$ сюръективно (7.4), а $\Gal(\Om'/F)\cong\varlim\Gal(E_k/F)$ (7.26a). Каждое отображение
$\Gal(E_{k+1}/F)\to\Gal(E_k/F)$ сюръективно с ядром порядка $[E_{k+1}:E_k]\ge2$, поэтому согласованное семейство строится
выбором на каждом шаге одного из не менее чем двух прообразов: таких семейств не меньше $2^{\aleph_0}$.
\end{remark}

\begin{proposition}[7.27]
Каждая проконечная группа $G$ является группой Галуа некоторого расширения Галуа.
\end{proposition}

\begin{proof}
Пусть $S=\bigsqcup_HG/H$~--- дизъюнктное объединение множеств левых смежных классов по всем открытым подгруппам $H$. Группа
$G$ действует на каждом $G/H$ левыми сдвигами, а значит, и на $S$.
\begin{itemize}[leftmargin=*]
\item \emph{Действие точное.} Если $g$ действует тривиально, то $gH=H$, т.\,е.\ $g\in H$, для каждой открытой $H$; а
пересечение открытых подгрупп проконечной группы тривиально (7.25, 7.25a). Здесь важно взять \emph{все} $G/H$: при одном $G/H$
ядро действия было бы $\bigcap_ggHg^{-1}\ne\{e\}$ (если $G$ бесконечна), и в итоге получилась бы лишь конечная группа
$G/\bigcap_ggHg^{-1}$ (см.\ 7.10a(a)).
\item \emph{Стабилизаторы открыты:} стабилизатор $xH\in G/H$ равен $xHx^{-1}$. (Регулярное действие $G$ на самой себе не
годится: стабилизаторы тривиальны, а $\{e\}$ не открыта в бесконечной проконечной группе.)
\end{itemize}
Пусть $k$~--- любое поле, $k[S]$~--- кольцо многочленов от переменных, занумерованных элементами $S$, и $E=k(S)$~---
его поле частных. Элемент $g\in G$ действует на $E$ подстановкой переменных:
$g\cdot P(s_1,\dots,s_n)=P(gs_1,\dots,gs_n)$; это автоморфизм $k[S]$, продолжающийся на поле частных. Получаем
гомоморфизм $G\to\Sym(S)\hookrightarrow\Aut(k(S))$ (перестановки переменных~--- лишь малая часть автоморфизмов $k(S)$,
но нам нужна только инъективность), так что $G$~--- группа автоморфизмов $E$. Стабилизатор любого $f\in E$ открыт: $f$
зависит лишь от конечного числа переменных $s_1,\dots,s_n$, поэтому
$\Stab(f)\supset\Stab(s_1)\cap\dots\cap\Stab(s_n)$~--- открытая подгруппа (7.1a(e)). Группа $G$ компактна. По 7.10
$E/E^G$~--- расширение Галуа с группой $G$ (и с той же топологией).
\end{proof}

\begin{remark}
Поле $F=k(S)^G$ здесь огромно и никак не контролируется; утверждение не говорит ничего о группах Галуа над заданным полем
(например, над $\Q$), и не каждая проконечная группа~--- \emph{абсолютная} группа Галуа (7.11). Для конечной $G$ можно
взять $E=k(x_g\mid g\in G)$ с действием перестановками; вопрос о том, чисто ли трансцендентно поле $E^G$ над $k$, известен
как \emph{проблема Нётер} (в общем случае ответ отрицательный).
\end{remark}
